%; whizzy paragraph -pdf xpdf -latex ./whizzypdfptex.sh
%; whizzy-paragraph "^\\\\begin{frame}"
% latex beamer presentation.
% platex, latex-beamer でコンパイルすることを想定。 

%     Tokyo Debian Meeting resources
%     Copyright (C) 2009 Junichi Uekawa

%     This program is free software; you can redistribute it and/or modify
%     it under the terms of the GNU General Public License as published by
%     the Free Software Foundation; either version 2 of the License, or
%     (at your option) any later version.

%     This program is distributed in the hope that it will be useful,
%     but WITHOUT ANY WARRANTY; without even the implied warranty of
%     MERCHANTABILITY or FITNESS FOR A PARTICULAR PURPOSE.  See the
%     GNU General Public License for more details.

%     You should have received a copy of the GNU General Public License
%     along with this program; if not, write to the Free Software
%     Foundation, Inc., 51 Franklin St, Fifth Floor, Boston, MA  02110-1301 USA

\documentclass[cjk,dvipdfmx,12pt]{beamer}
\usetheme{Tokyo}
\usepackage{monthlypresentation}

%  preview (shell-command (concat "evince " (replace-regexp-in-string "tex$" "pdf"(buffer-file-name)) "&"))
%  presentation (shell-command (concat "xpdf -fullscreen " (replace-regexp-in-string "tex$" "pdf"(buffer-file-name)) "&"))
%  presentation (shell-command (concat "evince " (replace-regexp-in-string "tex$" "pdf"(buffer-file-name)) "&"))

%http://www.naney.org/diki/dk/hyperref.html
%日本語EUC系環境の時
\AtBeginDvi{\special{pdf:tounicode EUC-UCS2}}
%シフトJIS系環境の時
%\AtBeginDvi{\special{pdf:tounicode 90ms-RKSJ-UCS2}}

\title{東京エリア Debian 勉強会}
\subtitle{資料}
\author{上川 純一 dancer@debian.or.jp\\IRC nick: dancerj}
\date{2009年5月16日}
\logo{\includegraphics[width=8cm]{image200607/openlogo-light.eps}}

\begin{document}

\frame{\titlepage{}}

\emtext{設営準備にご協力ください}

\begin{frame}{今日の参加の目標}
\begin{itemize}
 \item 上川: erlang パッケージの作成がしたいです。
 \item あけど: DDTSS の使い方をきちんと身につけたい。emacs/latex使えるようにな
       りたい。
 \item 山本: DDTSSの使い方を学びたい。
 \item 小室: Debianのリハビリ
 \item 日比野: ocaml 関係のパッケージの翻訳してみたい。
 \item でん: 最近 Android と言う人が多いのでAndroid SDK についておぼえてきた
       い。
 \item きたはら: Androidのあのじも聞かないけど興味があります。
 \item 前田: 4月から企画しているのでアンドロイドがどうこうというので興味
       あります。
 \item 藤澤: erlangの深淵にふれたい。
\end{itemize}
\end{frame}

\begin{frame}{今日の参加の目標}
\begin{itemize}
 \item やまだたくま: Androidの話がききたい。
 \item 吉田@板橋: twitter で話題の DDTSS という謎のキーワードを学びたい。
 \item masaka: xxxx
 \item やまねひでき: xxxx
\end{itemize}
\end{frame}

\section{}
\begin{frame}
 \frametitle{Agenda}
\begin{minipage}[t]{0.45\hsize}
  \begin{itemize}
  \item 注意事項
	\begin{itemize}
	 \item ごみは持ち帰りましょう。
	 \item となりで三味線をひいているので、おおきな声はだすな。
	\end{itemize}
  \item 最近あったDebian関連のイベント報告
	\begin{itemize}
	 \item 前回の勉強会
	\end{itemize}
 \end{itemize}
\end{minipage} 
\begin{minipage}[t]{0.45\hsize}
 \begin{itemize}
  \item Erlang
  \item MC-MPIパッケージの道
  \item Andoidネタ
  \item DDTSS翻訳プロセスワークショップ
 \end{itemize}
\end{minipage}
\end{frame}

\section{最近}

\begin{frame}
 \frametitle{2009年4月}
\begin{minipage}[t]{0.45\hsize}
  \begin{itemize}
  \item 注意事項
	\begin{itemize}
	 \item 飲食禁止
	 \item 政治/宗教/営利活動禁止
	\end{itemize}
  \item 最近あったDebian関連のイベント報告
	\begin{itemize}
	 \item 前回の勉強会
	 \item OSC
	 \item Ubuntu オフラインミーティング＠秋葉原
	 \item Linux Consortium 10 years event
	 \item Hack Cafe
	\end{itemize}
 \end{itemize}
\end{minipage} 
\begin{minipage}[t]{0.45\hsize}
 \begin{itemize}
  \item Debian Java policy について
  \item Debian 開発のワークフローについて語る会
  \item ocaml勉強はじめました
 \end{itemize}
\end{minipage}
\end{frame}

\begin{frame}{Hack Cafe}

毎週水曜日、週に一回東京のどっかのカフェでハック。

\end{frame}

\section{DWN quiz}
\begin{frame}{Debian 常識クイズ}

Debian の常識、もちろん知ってますよね?
知らないなんて恥ずかしくて、知らないとは言えないあんなことやこんなこと、
みんなで確認してみましょう。

今回の出題範囲は\url{debian-devel-announce@lists.debian.org} に投稿された
内容とDebian Project Newsからです。

\end{frame}

\subsection{問題}
%; whizzy-master ../debianmeetingresume200905.tex
% $B0J>e$N@_Dj$r$7$F$$$k$?$a!"$3$N%U%!%$%k$G(B M-x whizzytex $B$9$k$H!"(Bwhizzytex$B$,MxMQ$G$-$^$9!#(B
%
% $B$A$J$_$K!"%/%$%:$OJL%V%i%s%A$G:n@.$7!"$N$A$K%^!<%8$7$^$9!#5U$K%^!<%8$7(B
% $B$J$$$h$&$K$7$^$7$g$&!#(B
% (shell-command "git checkout quiz-prepare")

\santaku
{}
{}
{}
{}
{}

\begin{frame}{2009年計画}

{\scriptsize
 \begin{enumerate}
  \item 新年の企画 (アンサンブル荻窪開催)
  \item OSC Tokyo
  \item VAIO P インストール記録、
	カーネル読書会　ディストリビューション大集合(小林さん)(東京大学?)
  \item Git Handson (岩松)(あんさんぶる荻窪?)
  \item 家Debianサーバ vs 職場のネットワーク(千代田区都立図書館?\footnote{\url{http://www.library.chiyoda.tokyo.jp/}})
  \item Asterisk (東京大学?)
  \item スペインにて開催
  \item Debconf報告会
  \item OSC Fall?
  \item udev + HAL
  \item 3D graphics 開発 
  \item Debian サーバ+VMware + 各種OS、
	他の仮想化ツール(vserver etc.)、
	忘年会
 \end{enumerate}
}
\end{frame}

\emtext{MC MPI パッケージ ITP するぞ}

\emtext{Erlang を使ってみる}

\emtext{Android SDKをDebianで使ってみた}

\begin{frame}{Androidとは}
\begin{itemize}
 \item xxx
\end{itemize}
\end{frame}

\begin{frame}{Android SDK インストール}
Debian GNU/Linux でのインストールは結構簡単
\begin{itemize}
 \item eclipse 3.4 のパッケージがまだできていないので、自分でインストー
       ルする
 \item 必要な依存関係のパッケージをインストールする
\end{itemize}
\end{frame}


\begin{frame}{SDKを使ってみた}

\end{frame}



\begin{frame}{事前課題}

DDTSS
でいくつか(2個以上)Debianパッケージの説明文を翻訳してみて、2個レビューし
てみてください。その上で次の内容について説明してください:

\begin{enumerate}
 \item 適用した主要な方法
 \item 発見した課題
 \item 提案する理想像(ツールとか)、共有したい情報
\end{enumerate}
\end{frame}

{\footnotesize
%; whizzy-master ../debianmeetingresume200905.tex
% $B0J>e$N@_Dj$r$7$F$$$k$?$a!"$3$N%U%!%$%k$G(B M-x whizzytex $B$9$k$H!"(Bwhizzytex$B$,MxMQ$G$-$^$9!#(B

\begin{prework}{$B>e@n=c0l(B}

\preworksection{$BE,MQ$7$?<gMW$JJ}K!(B}

DDTSS $B$r%V%i%&%6$GD/$a$F!"K]Lu$r%l%S%e!<$7$F$_$^$7$?!#(B
$B%Q%C%1!<%8$N(Bdescription$B$@$1$GJ,$+$i$J$$ItJ,$K$D$$$F$O%P%0%l%]!<%H$r$7$^(B
$B$7$?!#(B
$BITL@E@$O(B debian-doc@jp $B%a!<%j%s%0%j%9%H$K<ALd$H$7$FEj9F$7$^$7$?!#(B

\preworksection{$BH/8+$7$?2]Bj(B}

$B%l%S%e!<$r$7$FJ,$+$C$?$3$H$G$9$,!"K]Lu$N:n6H$GC18l$NK]Lu$^$G$O$G$-$F$$$k(B
$B$N$G$9$,!"J8>OA4BN$H$7$F78$j<u$1$,$*$+$7$/!"0UL#$,DL$C$F$$$J$$$b$N$,$"$j$^$7$?!#(B
$B$^$?!"1Q8l$N6gFIE@$r$=$N$^$^F|K\8l$N6gFIE@$KCV$-49$($F$*$j!"$=$N$^$^$G$OF|K\8l$H$7$F(B
$BJ8>O$,D9$9$.$k$b$N$b$"$j$^$7$?!#K]Lu:n6H$rD>Lu:n6H$H$9$k$HFI$_$K$/$$(B
Description$B$,$G$-$"$,$C$F$7$^$&$H;W$o$l$^$9!#(B

\preworksection{$BDs0F$9$kM}A[A|(B($B%D!<%k$H$+(B)$B!"6&M-$7$?$$>pJs(B}

\begin{itemize}
 \item Description $B$KBP$7$F$9$G$K%P%0%l%]!<%H$,Ej9F$5$l$F$$$k$+$I$&$+$N(B
       $B%A%'%C%/$9$k%D!<%k!#(B

 \item $BMQ8l=8$K$9$G$K7G:\$5$l$F$$$kMQ8l$r4JC1$K%&%'%V%V%i%&%6$+$i%A%'%C%/(B
       $B$G$-$k(Bgreasemonkey $B!#(B

 \item DDTSS $B$G9T$C$?%l%S%e!<!&%3%a%s%H$J$I$,(B debian-doc@jp $B%a!<%j%s%0%j(B
       $B%9%H$KH?1G$9$k$3$H!#(B
\end{itemize}

\end{prework}

\begin{prework}{$B$^$($@$3$&$X$$(B}
\preworksection{$BE,MQ$7$?<gMW$JJ}K!(B}
$B;vA02]Bj$^$G<j$,2s$i$:$8$^$$!#$@$1$@$H7]$,$J$$$N$G!"8=:_M'?M$H?J$a$F$$$k(B
 CouchDB$B$N(BWeb$B%5%$%H$NK]Lu$K$D$$$FOC$7$^$9!#(B
\begin{enumerate}
 \item CouchDB$B$N(BWeb$B%5%$%H$r(BGoogle Docs$B$G$^$k$4$H<h$j9~$`!##1%Z!<%8$K$D$-!"(B
       $BK]Lu<T$O8BDj!#B>$O::FI$7!"0lJ8$:$DK]Lu$9$k!#D>Lu$G$O$J$/!"F|K\8l(B
       $B$H$7$F$o$+$j$d$9$$J8>O$r=E;k!#(B
 \item Web$B%5%$%H!"(BWiki$B$r$"$kDxEY$NHfN($^$GK]Lu$r?J$a$k!#(B
 \item CouchDB$B$N3+H/<T8~$1$N(BML$B$K8x3+$7$?$$;]$rAjCL$9$k!#(B
 \item Debian$B$N(BCouchDB$B$N%Q%C%1!<%8%a%s%F%J$G$b$"$k!"3+H/<T$+$i(Bpo4a$B$r;H$&(B
       $B$3$H!"B>$K$b$$$m$$$m2]Bj!J%7%9%F%`4IM}!"%"%C%W%m!<%I$NJ}K!!"K]Lu(B
       $B$NDDIe2=$J$I!K$,$"$k$1$I$=$l$r9MN8$7$J$$$H$$$1$J$$$h$H!"%"%I%P%$(B
       $B%9$r$b$i$&!#(B
 \item PO$B7A<0$G%Q%C%AEj$2$FH?1G<+BN$O$*4j$$$7!"99?7$O$J$s$i$+$N<jCJ$G3N(B
       $BG'$7$F?o;~K]Lu$7$F$$$/;]$rEA$($k!#(B
 \item $B:#8e$NJ}?K$r$I$&$9$k$+$rM'?M$HAjCL$7$F!":#%3%3!#(B
\end{enumerate}

\preworksection{$BH/8+$7$?2]Bj(B}
Web$B%5%$%H$O(BXHTML$B$K$J$C$F$$$J$$$N$G!"$^$:JQ49$9$k$H$3$m$+$i;O$a$J$$$H$$$1(B
 $B$J$$$3$H$H!"(BWiki$B$O99?7$,IQHK$K$"$j$9$.$@$+$i!"BP>]$+$i30$7$?J}$,NI$$$h(B
 $B$M!"$HM'?M$HAjCL$7$?$H$3$m!#K]Lu$,$=$b$=$b$d$j$?$$$3$H$G$O$J$+$C$?$N$G!#(B

\preworksection{$BDs0F$9$kM}A[A|(B($B%D!<%k$H$+(B)$B!"6&M-$7$?$$>pJs(B}
$B:G=i$NK]Lu$@$19T$C$?$i!"8e$O<jN%$l$9$k!JC/$+$K0z$-7Q$0!K$N$,M}A[$G$O$"$k(B
 $B$b$N$N!"6(NO<T$rA}$d$5$J$$$H$$$1$J$$!#$=$b$=$b(BDebian$B%Q%C%1!<%8$K$J$C$F(B
 $B$$$k$N$G!"%=!<%9%Q%C%1!<%8$K4^$^$l$k%I%-%e%a%s%H$J$i!"(BDebian-doc$B$GK]Lu(B
 $B$7$F!"$=$C$A7PM3$G%"%C%W%9%H%j!<%`$KH?1G$7$F$b$i$&!"$H$$$&$N$b<j$@$J$H(B
 $B$$$&$3$H$r8!F$Cf!#(B
\end{prework}

\begin{prework}{$B$"$1$I(B}

$B@bL@J8$rK]Lu$7$?%Q%C%1!<%8(B
libvolpack1	$B%l%s%@%j%s%0%i%$%V%i%j(B
libvolpack1-dev	$B%l%s%@%j%s%0%i%$%V%i%j(B
nagios3		$B%7%9%F%`4F;k(B/$B4IM}%D!<%k(B
nagios3-doc	$B%7%9%F%`4F;k(B/$B4IM}%D!<%k$NJ8=q(B
cwcp		$B%b!<%k%9?.9fN}=,%=%U%H(B
dnswalk		DNS$B%A%'%C%/%D!<%k(B
slapd		LDAP$B%5!<%P(B
pingpong	$B%"%^%A%e%"L5@~MQ(B convers $B%5!<%P(B
sugar-presence-service	OLPC$BMQ%=%U%H(B
amap-align	DNA$B2r@O%D!<%k(B
loki		DNA$B2r@O%D!<%k(B
xastir		DNA$B2r@O%D!<%k(B
biosquid	DNA$B2r@O%D!<%k(B
boxshade	DNA$B2r@O%D!<%k(B
dialign		DNA$B2r@O%D!<%k(B
dialign-t	DNA$B2r@O%D!<%k(B
probcon		DNA$B2r@O%D!<%k(B

$B!A%l%S%e!<$3$3$+$i!A(B
$B%l%S%e!<(B 1
$B%Q%C%1!<%8(B: libvolpack1

Short description
$B86J8(B: fast volume rendering library
$BLuJ8(B: $B9bB.$JBgNL%l%s%@%j%s%0%i%$%V%i%j(B

Long description
$B86J8(B:
VolPack is a software library for fast, high-quality volume rendering with
this features:
 * Renders data sampled on a regular, three-dimensional grid.
 * Supports user-specified transfer functions for both opacity and color.
 * Provides a shading model with directional light sources, multiple material
   types with different reflective properties, depth cueing, and shadows.
 * Produces color (24 bits/pixel) or grayscale (8 bits/pixel) renderings,
   with or without an alpha channel.
 * Supports arbitrary affine view transformations.
 * Supports a flexible data format that allows an arbitrary C structure to be
   associated with each voxel.
$BLuJ8(B:
VolPack $B$O9bIJ<A$G9bB.$JBgNL$N%l%s%@%j%s%0$r9T$&%=%U%H%&%'%"%i%$%V%i%j$G!"(B
$B<!$NMM$JFCD'$r;}$C$F$$$^$9(B:
 * $B@55,2=$5$l$?%5%s%W%k%G!<%?$r(B 3 $B<!85%0%j%C%I$KI=<((B
 * $BITF)L@$H%+%i!<N>J}$N%f!<%6Dj5A$NJQ494X?t$r%5%]!<%H(B
 * $BJ}8~;XDj$5$l$?8w8;$K$h$k1"1FIU$-%b%G%k!"0c$&H?<MFC@-$r;}$C$?J#?t$NAG:`%?(B
   $B%$%W!"G;C8!"5Z$S1F$rDs6!(B
 * $B%"%k%U%!%A%c%M%k$NM-$j!"L5$7$N(B (24 $B%S%C%H(B/$B2hAG$N(B) $B%+%i!<$^$?$O(B (8 $B%S%C%H(B
   /$B2hAG$N(B) $B%0%l!<%9%1!<%9$G$NIA2h$N@8@.(B
 * $BG$0U$N;kE@$N%"%U%#%sJQ49$r%5%]!<%H(B
 * $BG$0U$N(B C $B8@8lE*$J9=B$$r%\%/%;%k(B (3 $B<!853J;R>e$N:G>.N)J}BN!"(B2 $B<!85$G$N(B
   $B%T%/%;%k(B) $B$X7k$S$D$1$k;v$,=PMh$k=@Fp$J%G!<%?%U%)!<%^%C%H$N%5%]!<%H(B

$B%l%S%e!<(B 2
$B%Q%C%1!<%8(B: amap-align

Short description
$B86J8(B: Protein multiple alignment by sequence annealing
$BLuJ8(B: $B%7!<%/%(%s%9%"%K!<%j%s%0$K$h$kCAGr<A$N%^%k%A%W%k%"%i%$%a%s%H(B ($BB?=EG[Ns@0Ns(B)

Long description
$B86J8(B:
AMAP is a command line tool to perform multiple alignment of peptidic
sequences. It utilizes posterior decoding, and a sequence-annealing
alignment, instead of the traditional progressive alignment method. It
is the only alignment program that allows to control the sensitivity /
specificity tradeoff.  It is based on the ProbCons source code, but
uses alignment metric accuracy and eliminates the consistency
transformation.
.
 Homepage: http://bio.math.berkeley.edu/amap/
$BLuJ8(B:
AMAP $B$O%Z%W%A%IG[Ns$N%^%k%A%W%k%"%i%$%a%s%H(B ($BB?=EG[Ns@0Ns(B) $B$r9T$&%3%^%s%I(B
$B%i%$%s%D!<%k$G$9!#;v8e2rFI$KMxMQ$5$l!"EAE}E*$J%W%m%0%l%C%7%V%"%i%$%a%s%H(B
$BK!$NBe$o$j$K%7!<%/%(%s%9%"%K!<%j%s%0$N@0Ns$r9T$$$^$9!#46EY(B \& $BFC0[@-$N(B
$B%H%l!<%I%*%U$r%3%s%H%m!<%k=PMh$kM#0l$N%"%i%$%a%s%H(B ($B@0Ns(B) $B%W%m%0%i%`$G$9!#(B
ProbCons $B$N%=!<%9%3!<%I$r85$K$7$F!"@:L)B,Dj%"%i%$%a%s%H(B ($B@0Ns(B) $B$H%3%s%7(B
$B%9%F%s%7!<%H%i%s%9%U%)!<%a!<%7%g%s(B ($B0l4S@-JQ2=(B) $B$N=|5n$rMxMQ$7$F$$$^$9!#(B
.
 Homepage: http://bio.math.berkeley.edu/amap/

$B!A%l%S%e!<$3$3$^$G!A(B

1$B!#E,MQ$7$?<gMW$JJ}K!(B
$B<jF~NO$H%M%C%H>e$N1QOB<-E5!"I,MW$K1~$8$F(B Google$B!"(Bwikipedia$B!"$=$NB>@lLg%5%$%H$r;2>H(B

2$B!#H/8+$7$?2]Bj(B
DNA$B2r@O%D!<%k$,B?$/$"$j!"@lLgMQ8l$,$^$A$^$A$G(B
$B$=$l$NE}0l<c$7$/$O6&DL2=$,I,MW$@$H;W$$$^$7$?!#(B

3$B!#Ds0F$9$kM}A[A|(B($B%D!<%k$H$+(B)
$B@lLgMQ8l$K$D$$$F$O<-=q%Z!<%8$rMQ0U$9$k$+!"4p=`$K$9$k@lLgMQ8l$N%5%$%H$r7h$a$k$N$,NI$$$+$H;W$$$^$9!#(B

$B!&6&M-$7$?$$>pJs$r@bL@(B
Debian $B%Q%C%1!<%8$K$O?tB?$/$N(B DNA $B2r@O%D!<%k$,$"$C$F!"@lLgMQ8l$NK]Lu$K$"$?$C$FCm0U?<$/:n6H$9$kI,MW$,$"$k$H;W$$$^$9!#(B
DNA $B2r@O$K$D$$$FM=HwCN<1$H$7$F2<5-%5%$%H$,;29M$K$J$k$+$H;W$$$^$9!#(B
$B%?%s%Q%/<A!?3K;@$N%"%i%$%a%s%H2r@O(B 
http://www.icot.or.jp/ARCHIVE/Museum/SOFTWARE/GIP/gene\_alignment.html
$B"(J8;z%3!<%I$,(B ISO-2022-JP $B$J$N$K%(%s%3!<%I>pJs$O(B Shift-JIS $B$K$J$C$F$$$^$9!#(B

Debian JP $BJ8=q:n@.$N;X?K(B $B$K$D$$$F(B DDTSS for ja $B$K%j%s%/$,$"$C$?J}$,NI$$$+$H;W$$$^$9!#(B

\end{prework}

\begin{prework}{$BF|HfLn(B $B7<(B}
DDTSS$B$G(Bocaml$B$H(Blibperl4caml-ocaml$B$N(BDescription$B$rK]Lu$7$F$_$^$7$?!#(B
$B%l%S%e!<$O(Bnasm$B$H(Blibobjc2$B$r$d$C$F$_$^$7$?!#(B

\preworksection{$BE,MQ$7$?<gMW$JJ}K!(B}

DDTSS$B$N(BWeb$B%$%s%?!<%U%'!<%9$+$iK]Lu$H%l%S%e!<$r9T$J$C$F$_$^$7$?!#(B

\preworksection{$BH/8+$7$?2]Bj(B}

$B<+J,$,K]Lu$r9T$J$C$?FbMF$,99?7$5$l$?$H$-$d(B
$B%3%a%s%H$,IU$$$?$H$-$K!":F$S8+$K$$$+$J$$$H99?7$rCN$k$3$H$,(B
$B$G$-$J$$$N$,5$$K$J$j$^$7$?!#(B

\preworksection{$BDs0F$9$kM}A[A|(B($B%D!<%k$H$+(B)$B!"6&M-$7$?$$>pJs(B}

$BK]Lu$7$??M$d%l%S%e!<$7$??M$K99?7$rDLCN$9$k5!G=$H(B
$B$=$l$r%3%s%H%m!<%k$9$k5!G=$,M_$7$$$G$9!#(B
\end{prework}

\begin{prework}{$B>.@n(B $B?-0lO:(B}
 $B0J2<!$$d$C$F$_$?46A[$J$I!%(B
 \preworksection{$BEPO?(B}
 $B$^$:EPO?;~$N(B Alias $B$H8@$&$N$,D>46$G$o$+$j$^$;$s$G$7$?!%(B
 \preworksection{$BK]Lu(B}
 $BK]Lu$O0J2<$N#2$D$N%Q%C%1!<%8$K$D$$$F$d$C$F$_$^$7$?!%(B
 \begin{itemize}
  \item wgerman-medical
  \item libbio-ruby
 \end{itemize}
 $B;W$C$?$N$O!$$$$-$J$jA4$F$rLu$5$J$1$l$P$J$i$J$$$N$O?I$$$s$8$c$J$$$+$H8@(B
 $B$&$3$H!%>/$7$:$DK]Lu$G$-$?$j!$$b$7$/$OESCf7k2L$rJ]B8$G$-$k$?J}$,$$$$$s(B
 $B$8$c$J$$$+$H;W$$$^$7$?!%(B

 $B$"$H$O!$:G=i$K$I$&2?$r$9$l$P$$$$$N$+$,$o$+$j$K$/$H8@$&$3$H!%(B
 $B$o$+$C$?Ev$?$jA0$J$N$+$bCN$l$^$;$s$,!$I}9-$/6(NO<T$r5a$a$k$N$J$i!$2?$+(B
 $B$7$i>pJs(B(HowTo$B$H$+!)(B)$B$"$C$?J}$,$$$$5$$,$7$^$9!%(B
 \preworksection{$B%l%S%e!<(B}
 $B$$$-$J$j(B 'Accept as is' $B$H=q$+$l$k$N$,I]$+$C$?!%$A$c$s$H!$(B
 ``'Accept as is' means you agree with this translation.''
 $B$H=q$+$l$F$O$$$k$s$G$9$,!$$I$&$K$b!V$($C!$$$$$$N$+!)!W$H8@$&;W$$$,H4$1(B
 $B$-$l$^$;$s$G$7$?!%$N$G!$%l%S%e!<$O8+$?$@$1$G$9!%(B

 $B$G!$(Bmultimail $B$r8+$?$H$-$N46A[$G$9$,!$(Bdiff $B$r$b$&>/$7$o$+$j$d$9$/I=<($G(B
 $B$-$J$$$+$H;W$$$^$7$?!%(B
 $B$5$9$,$KA08e$K$"$k$H!$F,$NCf$G$$$m$$$mAH$_N)$F$J$$$HFI$a$J$$$N$G!$(B
 $B>e2<$K0[$J$kItJ,$@$1J8;zD9$rJ;$;$FI=<($G$-$?$j$9$k$H!$$o$+$j$d$9$$$+$J(B
 $B$H!%(B

 $B!V;W$C$?!W$P$+$j$N46A[$K$J$C$F$7$^$$$^$7$?$,!$$3$s$J46$8$G$7$?!%(B
\end{prework}

\begin{prework}{$B$d$^$@$?$/$^(B}

\preworksection{$BE,MQ$7$?<gMW$JJ}K!(B}

DDTSS $B$h$j86J8$HLuJ8$rF~<j$7$F!"K]Lu;Y1g%D!<%k(B OmegaT $B$^$?$O%F%-%9%H(B
$B%(%G%#%?$GK]Lu$7$^$7$?!#J#?t$N%*%s%i%$%sHG$N<-=q!"%Z!<%Q!<HG$N<-=q!"K]Lu(B
$B%=%U%H!"BP>]%=%U%H%&%'%"!"(BWeb $B8!:w!"=q@R$ND4::$rJ;MQ$7$^$7$?!#(B
$B%l%S%e!<0MMj$r(B debian-doc $B%a!<%j%s%0%j%9%H$KEj9F$7$^$7$?!#(B
$B>pJsDs6!0MMj$r(B YLUG $B%a!<%j%s%0%j%9%H$KEj9F$7$^$7$?!#(B

\preworksection{$BH/8+$7$?2]Bj(B}

Deiban JP $B4XO"$NK]Lu$N%7%9%F%`$H1?MQ$K$D$$$F$N>pJs<}=8$NFq0WEY$,9b$$(B
$B>uBV$G$7$?!#%I%-%e%a%s%H$,@0Hw$5$l$F$$$J$$>pJs$d!"?tG/4V%"%C%W%G!<%H$5$l$F(B
$B$$$J$$%Z!<%8$,$"$j$^$7$?!#(B

$BMW5a$5$l$kLu$NIJ<A$,ITL@$J$?$a!":n6H;~4V$d:n6H<j=g$N8+@Q$j$,:$Fq$G$7$?!#(B

$BLu$NIJ<A0BDj$KI,MW$J<9I.4p=`$N@0Hw>u67$,IT==J,$G$7$?!#(B

\preworksection{$BDs0F$9$kM}A[A|(B($B%D!<%k$H$+(B)$B!"6&M-$7$?$$>pJs(B}

\begin{itemize}
 \item $B:n6H<j=g%Z!<%8$N%"%C%W%G!<%H!#(B

 \item $B<9I.4p=`$HMQ8l=8$N%"%C%W%G!<%H!#(B

 \item $B<9I.4p=`E,9g%F%9%H%D!<%k$dK]Lu;Y1g%D!<%kMQ%i%$%V%i%j$NDs6!$,$G$-(B
       $B$k$H$&$l$7$$$G$7$g$&!#(B
\end{itemize}

\end{prework}

\begin{prework}{$B$d$^$M$R$G$-(B}

\preworksection{$BE,MQ$7$?<gMW$JJ}K!(B}

DDTSS $B$r%V%i%&%6$GD/$a$F!"K]Lu!u%l%S%e!<!D$7$?$h$&$J5-21$,Hy$+$K$"$j$^$9!#(B

\preworksection{$BH/8+$7$?2]Bj(B}

$B0l?M$@$H<d$7$9$.$k!#$H$$$&$N$O$5$F$*$-!"A4BNA|$H8=>u$NGD0.$,Fq$7$/!"0lBN$I$3$+$i$3$N;3$r(B
$BEP$l$P$$$$$N$+$h$/J,$+$i$J$$>uBV$@$C$?5-21$,!#(B

\preworksection{$BDs0F$9$kM}A[A|(B($B%D!<%k$H$+(B)$B!"6&M-$7$?$$>pJs(B}

\begin{itemize}
 \item $B8+$d$9$$%Z!<%89=@.$,M_$7$$!#AGKQ2a$.$F2?$,2?$d$iJ,$+$i$s$N$O4*J[!#$;$a$F(Bpo-debconf$B$N%Z!<%8(B(http://www.debian.org/international/l10n/po-debconf/ja)$B%l%Y%k$N$OM_$7$$!#(B
 \item $B!V0l%f!<%6$,FI$s$GM}2r$G$-$k!WJ8>O$r?4$,$1$k$3$H!#Ev$?$jA0$J$s$G$9$,!"Fq$7$$$N$G2~$a$F!J5U$K8@$&$H!"::FI$O%f!<%6$K$7$F$b$i$&;EAH$_$,M_$7$$$H$3$m$G$9!K!#(B
\end{itemize}

\end{prework}

}

\emtext{DDTSS ワークショップ}


\begin{frame}{次回勉強会のテーマ}
\begin{itemize}
 \item 会場は市ヶ谷
 \item MC MPI ネタ続き?
 \item DDTSSの続き?
\end{itemize}
\end{frame}

\begin{frame}{Debconf 2009}
\begin{itemize}
 \item スペイン参加
 \item 準備は?
\end{itemize}
\end{frame}


\begin{frame}{宴会場所}
\begin{itemize}
 \item 宴会場所\\
       本日の宴会は「のみくい処 仲々」です。\\
       21:00 開始です。
 \item 片付け\\
       部屋を片付けるのにご協力ください。
\end{itemize}
\end{frame}

\end{document}

;;; Local Variables: ***
;;; outline-regexp: "\\([ 	]*\\\\\\(documentstyle\\|documentclass\\|emtext\\|section\\|begin{frame}\\)\\*?[ 	]*[[{]\\|[]+\\)" ***
;;; End: ***
